% !TeX root = ../../main.tex
\section{Medie e Medie Condizionali}

Qui si applica lo stesso identico ragionamento visto per le probabilità, ma al \textbf{Valore Atteso} (la Media). Questa è una delle formule più usate al mondo, spesso chiamata \textbf{Teorema della Media Totale}.

\paragraph{(Esempio della Scuola)}
Immagina di voler calcolare l'altezza media di \textbf{tutti} gli studenti di una scuola ($\mathbb{E}[X]$).
\\È difficile misurarli tutti insieme. È più facile fare così:

\begin{enumerate}
    \item Vai nella \textbf{Classe 1} ($E_1$), calcoli l'altezza media lì dentro ($\mathbb{E}[X \mid E_1]$) e la pesi per quanti alunni ha la classe.
    \item Vai nella \textbf{Classe 2} ($E_2$), calcoli la media lì ($\mathbb{E}[X \mid E_2]$) e pesi.
    \item Alla fine sommi i risultati pesati.
\end{enumerate}

\paragraph{La Formula}
Matematicamente, il teorema si esprime come:

\begin{equation}
\boxed{\mathbb{E}[X] = \sum_{m=1}^M \mathbb{P}(E_m) \cdot \mathbb{E}[X \mid E_m]}
\end{equation}

Dove:
\begin{itemize}
    \item $\mathbb{P}(E_m)$: Il "peso" dello scenario (es. quanto è grande la classe).
    \item $\mathbb{E}[X \mid E_m]$: La media calcolata \textit{solo} dentro quello scenario.
\end{itemize}

\subsection{Derivazione Matematica}

Partendo dalla definizione di valore atteso:

\[
\mathbb{E}[X] = \sum_{x \in \mathcal{X}} x p_X(x)
\]

Sostituiamo $p_X(x)$ con la regola della probabilità totale vista prima e scambiamo l'ordine delle sommatorie:

\begin{equation}
\begin{aligned}
\mathbb{E}[X] &= \sum_{x \in \mathcal{X}} x \sum_{m=1}^M p_{X|E_m}(x)\mathbb{P}(E_m) \\
&= \sum_{m=1}^M \mathbb{P}(E_m) \underbrace{\sum_{x \in \mathcal{X}} x p_{X|E_m}(x)}_{\mathbb{E}[X \mid E_m]}
\end{aligned}
\end{equation}

\paragraph{Esempio finale (Binomiale)}
Se applichi questa formula agli esempi precedenti ($X \sim \mathcal{B}(16, p)$), puoi ritrovare la media generale scomponendola nei due casi:

\[
\mathbb{E}[X] = \mathbb{P}(X \le 4)\mathbb{E}[X \mid X \le 4] + \mathbb{P}(X > 4)\mathbb{E}[X \mid X > 4]
\]

Oppure in egual modo:
\[
\mathbb{P}(X \in \{2,3,4,5,6\}) \mathbb{E}[X \mid X \in \{2,3,4,5,6\}] + \mathbb{P}(X \notin \{2,3,4,5,6\}) \mathbb{E}[X \mid X \notin \{2,3,4,5,6\}]
\]

Significa: La media generale dei successi (che sappiamo essere $n \cdot p = 16/3 \approx 5.33$) è uguale alla media dei successi "quando va male" (pesata per la probabilità che vada male) PIÙ la media dei successi "quando va bene" (pesata per la probabilità che vada bene).

